% ==================== 3.4 月面光照与温度环境建模 ====================
\section{月面光照与温度环境建模}
\label{sec:light_temp}

月面光照与温度环境是影响巡视器能源供给和热控设计的关键因素。本节从太阳运动规律、光照场建模、温度场建模三个方面展开研究。

\subsection{太阳运动规律分析}
\label{subsec:sun_motion}

月球绕地球公转的同时自转，形成独特的太阳视运动规律。

\textbf{（1）太阳高度角}

太阳高度角定义为太阳光线与当地水平面的夹角，计算公式为：
\begin{equation}
    \sin h = \sin\phi\sin\delta + \cos\phi\cos\delta\cos\omega
    \label{eq:sun_altitude}
\end{equation}
其中，$\phi$为纬度，$\delta$为太阳赤纬，$\omega$为时角。

对于月球，太阳赤纬变化范围为$\pm 1.5\degree$（黄白交角），时角周期为月球自转周期（约27.3天）。

\textbf{（2）太阳方位角}

太阳方位角定义为太阳光线在水平面上的投影与正北方向的夹角：
\begin{equation}
    \cos A = \frac{\sin h \sin\phi - \sin\delta}{\cos h \cos\phi}
    \label{eq:sun_azimuth}
\end{equation}

\textbf{（3）月昼月夜周期}

月球自转周期约为27.3天，因此月昼和月夜各约13.65天。在极区，由于太阳高度角常年较低，存在持续光照区和永久阴影区。

\subsection{光照场建模}
\label{subsec:illumination}

光照场建模的目标是计算月面任意位置、任意时刻的太阳辐照度和阴影分布。

\textbf{（1）太阳辐照度计算}

月面接收的太阳辐照度取决于太阳高度角和地形遮挡。太阳辐照度可表示为：
\begin{equation}
    I = I_0 \cdot \sin h \cdot (1 - S)
    \label{eq:irradiance}
\end{equation}
其中，$I_0$为太阳常数（约1361 W/m²），$S$为阴影因子（被遮挡时$S=1$，否则$S=0$）。

\textbf{（2）阴影计算}

采用光线追踪方法计算阴影分布。对于月面上任意点$P$，沿太阳光线方向追踪，若遇到地形障碍则$P$处于阴影中。

阴影判断的条件为：
\begin{equation}
    h_{obs} \geq h_P + d_{obs-P} \cdot \tan h
    \label{eq:shadow}
\end{equation}
其中，$h_{obs}$为障碍物高程，$h_P$为点$P$的高程，$d_{obs-P}$为障碍物到点$P$的距离，$h$为太阳高度角。

\textbf{（3）光照时序分析}

对于规划区域，分析各点的光照时间分布，识别长时间处于阴影或光照的区域，为巡视器能源规划提供依据。

\subsection{温度场建模}
\label{subsec:temperature}

温度场建模的目标是计算月面任意位置、任意时刻的温度分布。

\textbf{（1）热平衡方程}

月面温度由热平衡方程确定：
\begin{equation}
    \frac{\partial T}{\partial t} = \frac{1}{\rho c_p}\frac{\partial}{\partial z}\left(k\frac{\partial T}{\partial z}\right) + \frac{Q}{\rho c_p}
    \label{eq:heat_eq}
\end{equation}
其中，$T$为温度，$\rho$为密度，$c_p$为比热容，$k$为热导率，$Q$为热源项（太阳辐射）。

\textbf{（2）边界条件}

表面边界条件：
\begin{equation}
    -k\left.\frac{\partial T}{\partial z}\right|_{z=0} = \alpha I - \epsilon\sigma T^4
    \label{eq:surface_bc}
\end{equation}
其中，$\alpha$为吸收率（约0.9），$\epsilon$为发射率（约0.9），$\sigma$为Stefan-Boltzmann常数（$5.67\times 10^{-8}$ W/m²/K⁴）。

深层边界条件：
\begin{equation}
    T|_{z\to\infty} = T_{deep}
    \label{eq:deep_bc}
\end{equation}
其中，$T_{deep}$为深层温度（约250 K）。

\textbf{（3）数值求解}

采用有限差分方法求解热平衡方程。将月壤沿深度方向离散为多层，时间采用显式或隐式格式推进。

月面温度典型值为：白天最高约390 K（117\degree C），夜间最低约100 K（-173\degree C），平均约220 K（-53\degree C）。
